% Одностраничный отчёт о научной работе.
% Все поля, которые обычно приходится менять, собраны в этом блоке.
\documentclass[11pt,letterpaper]{article}

\usepackage[T2A]{fontenc}
\usepackage[utf8]{inputenc}
\usepackage[russian]{babel}
\usepackage{amsmath}
\usepackage{geometry}
\usepackage{microtype}

\geometry{
  top=20mm,
  bottom=17mm,
  left=25mm,
  right=25mm
}

\newcommand{\StudentName}{Буканов Иван Алексеевич}
\newcommand{\StudentShortName}{Буканов И.\,А.}
\newcommand{\StudentOrganization}{ФКИ МГУ}
\newcommand{\ProjectPlace}{Лаборатория внегалактической астрофизики и космологии САО РАН}
\newcommand{\WorkPeriod}{13 июля - 3 августа 2026 г.}
\newcommand{\ProjectTitle}{Прямая калибровка метода флуктуаций поверхностной яркости по данным JWST}
\newcommand{\SupervisorName}{Макаров Д.\,И.}
\newcommand{\SupervisorDate}{«03» августа 2026 г.}
\newcommand{\StudentDate}{«03» августа 2026 г.}

\setlength{\parindent}{1.25cm}
\setlength{\parskip}{0pt}
\pagestyle{plain}

\begin{document}

\begin{center}
  {\LARGE Отчёт о работе над научным проектом}
\end{center}

\vspace{4mm}

\noindent
Студент \StudentName, \StudentOrganization

\vspace{4mm}

В период \WorkPeriod{} я выполнял работу в \ProjectPlace{} под руководством
\SupervisorName{} по теме «\ProjectTitle». Целью работы было реализовать и
проверить метод измерения расстояний до галактик по флуктуациям поверхностной
яркости (Surface Brightness Fluctuations, SBF), используя изображения
JWST/NIRCam и независимую шкалу расстояний по вершине ветви красных гигантов
(TRGB).

Для анализа использованы публичные изображения 14 галактик ранних
морфологических типов из программы JWST GO-3055 в фильтрах F090W и F150W.
Был разработан воспроизводимый программный конвейер на языке Python. Он
включает оценку и вычитание фона, построение масок дефектных пикселей и
компактных источников, аппроксимацию распределения света галактики изофотами,
построение гладкой двумерной модели и анализ нормированных остатков в
пространстве пространственных частот. Амплитуда SBF определялась по спектру
мощности
\[
  P(k)=P_0E(k)+P_1,
\]
где $E(k)$ учитывает форму функции рассеяния точки и маску, а из $P_0$
вычитался вклад неразрешённых шаровых скоплений и фоновых галактик. Измерения
проводились в двух круговых кольцах $8.2$--$16.4$ и $16.4$--$32.8$ угл. сек,
как в работе Jensen et al. по проекту TRGB--SBF. Для каждой галактики
сохранялись промежуточные изображения, каталоги источников, параметры
аппроксимации и журнал выполнения.

Абсолютные флуктуационные звёздные величины в F150W были получены с
использованием индивидуальных расстояний TRGB из четвёртой работы проекта
TRGB--SBF. Учтено поглощение в Млечном Пути; отдельно проверены устойчивость к
выбору модели функции рассеяния точки, вклад неразрешённых источников и
структурные остатки модели галактики. Для полной выборки получен нуль-пункт
\[
  \overline{M}_{150}=-3.172\ \text{зв. вел.}
\]
с внутренним разбросом $0.084$ зв. вел. Проверка с последовательным
исключением одной галактики даёт RMS различий между SBF- и TRGB-расстояниями
$0.096$ зв. вел., что соответствует $4.4\%$ по расстоянию; медианное абсолютное
отклонение составляет $3.5\%$. Линейная зависимость от цвета
F090W--F150W статистически не улучшает предсказание, поэтому основной
результат сформулирован как цветонезависимая калибровка; цветовые модели
оставлены как проверка чувствительности.

По итогам работы подготовлены таблицы измерений и ошибок для всех 14 галактик,
диагностические графики, автоматический пакетный запуск и черновик статьи.
Дальнейшая работа связана с независимой проверкой калибровки на новых
наблюдениях JWST и с улучшением двумерного моделирования галактик со сложной
крупномасштабной структурой.

\vfill

\noindent
\begin{minipage}[t]{0.48\textwidth}
  \SupervisorDate

\vspace{11mm}

  \StudentDate
\end{minipage}%
\hfill
\begin{minipage}[t]{0.46\textwidth}
  Руководитель практики:

  \vspace{2mm}
  \rule{35mm}{0.4pt}\,/ \SupervisorName

  \vspace{8mm}

  Студент:

  \vspace{2mm}
  \rule{35mm}{0.4pt}\,/ \StudentShortName
\end{minipage}

\end{document}
